\documentclass[11pt]{article}
\usepackage[a4paper,margin=1in]{geometry}
\usepackage{amsmath,amssymb,booktabs,graphicx,hyperref,xcolor}
\title{Autonomous Diffusion: the FID-faithful risk measure\\
       {\large Why the KL/$L^2$ certificate is not an FID certificate,
       and the feature-space Wasserstein risk that is}}
\author{thermau5}
\date{\today}
\begin{document}
\maketitle

\paragraph{Status.} This is a \emph{theory + diagnostic} report, separate
from the locked-test report (\texttt{report\_locked.tex}). It does not
re-state the locked empirical wins (those stand on measured Clean-FID and
are unaffected). It addresses one foundational question surfaced during
Round~6: \textbf{which risk measure is faithful to FID?} We show, from
first principles and with a controlled diagnostic, that the KL/$L^2$
calibrated certificate of PDF~1 bounds NLL/pathwise-$L^2$ risk but
\emph{not} FID, identify the exact structural reason, and derive +
validate the FID-faithful risk: a Wasserstein distance in
Inception-feature space with terminal error propagation. This is the
PDF~2 Wasserstein certificate evaluated in feature space.

\section{The mismatch, stated precisely}
\label{sec:mismatch}

The PDF~1 certificate is
\[
Q_s(u) = Q_0 + \frac{c_{\mathrm{stat}}}{n}\!\int_0^1 \frac{a(r)}{\rho(r)}\,dr
       + c_{\mathrm{disc}}\!\int_0^1 \frac{d_s(r)}{m(r)^p}\,dr .
\]
Only the discretization term depends on the sampling grid $m$; $Q_0$
(irreducible floor) and the statistical term (frozen by the pretrained
net) are $m$-independent. So across grids the certificate predicts
$R(m)-R_{\mathrm{floor}} = c_{\mathrm{disc}}\!\int d_s/m^p$.

\paragraph{The certificate is correctly computed --- as a pathwise-$L^2$
bound.} On a panel of $K{=}5$ Heun grids spanning FID $13\to270$, the
discretization term (computed correctly, including the terminal gap
$\sigma_{\text{last}}\!\to\!0$) tracks the true per-sample pixel-$L^2$
global error almost exactly. There is no implementation bug in $Q$.

\paragraph{But it does not track FID.} Two grids with \emph{equal}
certificate value have $6\times$ different FID:

\begin{center}
\small
\begin{tabular}{l r r}
\toprule
grid & $\int d_s/m^p$ (pixel, terminal-incl) & Clean-FID \\
\midrule
\texttt{good\_pointwise} & 26.03 & \textbf{13.2} \\
\texttt{forced\_smin}    & 26.41 & \textbf{87.9} \\
\bottomrule
\end{tabular}
\end{center}

\noindent The certificate cannot distinguish them. Across the full panel
the rank order of $\int d_s/m^p$ disagrees with the FID rank order
(e.g.\ \texttt{cluster\_hi} has middling certificate but the worst FID).
This is not a bug; it is a \textbf{risk-measure} problem.

\section{What FID is}
\label{sec:fid}

FID is not a pixel distance and not a per-sample error. With
$\varphi:\mathbb R^{d}\to\mathbb R^{2048}$ the Inception-v3 feature map,
and $(\mu_g,\Sigma_g)$, $(\mu_d,\Sigma_d)$ the mean/covariance of features
of generated and data samples,
\[
\mathrm{FID}
= \lVert \mu_g-\mu_d\rVert^2
 + \operatorname{Tr}\!\Bigl(\Sigma_g+\Sigma_d-2(\Sigma_g^{1/2}\Sigma_d\Sigma_g^{1/2})^{1/2}\Bigr)
= W_2^2\!\Bigl(\mathcal N(\mu_g,\Sigma_g),\,\mathcal N(\mu_d,\Sigma_d)\Bigr).
\]
That is, FID \emph{is} the squared $2$-Wasserstein distance between two
Gaussians fit in \textbf{Inception-feature space}. Three structural
properties follow, each of which the pixel/KL certificate violates:

\begin{enumerate}
\item \textbf{Space.} FID lives in feature space $\varphi(x)$, not pixels.
\item \textbf{Propagation.} It depends on the \emph{final} sample's
features; a step error must be propagated through the remaining flow and
through $\varphi$, not measured locally.
\item \textbf{Distributional.} It is a mean-shift plus covariance
mismatch, not a sum of per-sample errors. Zero-mean errors that average
out cost little FID; small \emph{biased} errors cost a lot.
\end{enumerate}

\section{The FID-faithful risk (first principles)}
\label{sec:derivation}

Let $\delta_i$ be the local truncation error of step $i$ (at noise level
$\sigma_i$), with $\mathbb E\lVert\delta_i\rVert^2 \sim d_s(\sigma_i)\,h_i^{p+1}$.
It propagates to the final sample $x_0$ through the remaining flow
Jacobian $A_{\sigma_i\to 0}$, then to features through the Inception
Jacobian $J_\varphi$. Thus $\Delta\varphi = \sum_i J_\varphi\,A_{\sigma_i\to 0}\,\delta_i$
and, to leading order,
\[
\mathrm{FID}_{\mathrm{disc}}
\;\approx\; \bigl\lVert \mathbb E[\Delta\varphi]\bigr\rVert^2
            + \operatorname{Tr}\operatorname{Cov}(\Delta\varphi)
\;\approx\; \sum_i \mathbb E\bigl\lVert J_\varphi\,A_{\sigma_i\to 0}\,\delta_i\bigr\rVert^2 .
\]
Defining the \textbf{feature-amplified difficulty}
\[
\boxed{\;d_s^{\varphi}(\sigma) \;=\; d_s(\sigma)\,\cdot\, g(\sigma),
\qquad g(\sigma) := \bigl\lVert J_\varphi\,A_{\sigma\to 0}\bigr\rVert^2\;}
\]
the FID-faithful discretization risk and its optimizer are
\[
Q^{\varphi}_s(u)=Q_0+\frac{c_{\mathrm{stat}}}{n}\!\int\frac{a}{\rho}
+c_{\mathrm{disc}}\!\int\frac{d_s^{\varphi}(\sigma)}{m(\sigma)^p}\,d\sigma,
\qquad
m^{\star}_{\mathrm{FID}}(\sigma)\propto d_s^{\varphi}(\sigma)^{1/(p+1)} .
\]
The factor $g(\sigma)=\lVert J_\varphi A_{\sigma\to0}\rVert^2$ ---
``how much a unit perturbation at noise level $\sigma$ moves the final
sample's Inception features'' --- is the principled FID-sensitivity
weight. \textbf{The hand-tuned perceptual weight $w(\sigma)=\sigma^{-2}$
(the $k{=}2$ knob of \texttt{report\_locked.tex}) is a scalar power-law
proxy for $g(\sigma)$.}

\paragraph{This is the PDF~2 Wasserstein certificate.} PDF~2 gives
$m^{\star}_W(r)\propto\bigl(A_y(r)\,\tau_y(r)\bigr)^{2/(p+1)}$ with terminal
amplification $A_y$. Since FID is a feature-space $W_2$, the FID-faithful
risk \emph{is} the Wasserstein variant with the amplification taken in
Inception-feature space, $A_y \to \lVert J_\varphi A_{\sigma\to0}\rVert$.
PDF~1 (KL/NLL) is faithful to likelihood; PDF~2 (Wasserstein), in feature
space, is faithful to FID.

\section{Estimating $g(\sigma)$}
\label{sec:estimator}

$g(\sigma)$ is directly measurable by a finite-difference Jacobian-vector
product through flow$+$feature map. At noise level $\sigma$, perturb a
reference state $x_\sigma^{\mathrm{ref}}$ by $\epsilon v$ ($v$ a random
unit vector), integrate the probability-flow ODE accurately to $\sigma{=}0$,
and measure the induced feature change:
\[
\hat g(\sigma)
=\frac{1}{\epsilon^2}\,
\mathbb E_{v}\bigl\lVert \varphi\bigl(x_0(x_\sigma^{\mathrm{ref}}+\epsilon v)\bigr)
                         -\varphi\bigl(x_0(x_\sigma^{\mathrm{ref}})\bigr)\bigr\rVert^2 .
\]
No hand-tuned weight is needed. The measured profile on EDM CIFAR-10
(28 noise levels, finite-difference, Heun terminal integration):
\[
g(\sigma{=}40)\approx 0.016,\qquad
g(\sigma{=}1)\approx 9.8,\qquad
g(\sigma{=}0.05)\approx 2.0\times10^{3}.
\]
Feature sensitivity is enormous at low $\sigma$ and negligible at high
$\sigma$: a perturbation near the end of sampling maps almost directly to
a feature change, while a high-$\sigma$ perturbation is washed out by the
long remaining flow. (This is steeper than $\sigma^{-2}$, explaining why a
fixed $k$ both helps and mis-calibrates.)

\section{Validation}
\label{sec:validation}

\noindent\textit{Setup:} \textsc{fix} $\theta_{\mathrm{EDM}},\pi_{\mathrm{EDM}}$,
Heun core, $K{=}5$; \textsc{variable} = 5 grids spanning FID $13\to270$;
\textsc{select} = none. We compare the pixel certificate $\int d_s/m^p$
and the feature risk $\int d_s^{\varphi}/m^p$ (both terminal-gap-included)
against measured Clean-FID.

\begin{center}
\small
\begin{tabular}{l r r r}
\toprule
grid & pixel $\int d_s/m^p$ & \textbf{feature $\int d_s^{\varphi}/m^p$} & Clean-FID \\
\midrule
\texttt{good\_pointwise} & 26.03 & \textbf{1.46} & 13.2 \\
\texttt{refined\_lowQ}   & 12.97 & \textbf{0.82} & 26.7 \\
\texttt{karras\_rho7}    & 64.11 & \textbf{2.74} & 44.9 \\
\texttt{forced\_smin}    & 26.41 & \textbf{4.00} & 87.9 \\
\texttt{cluster\_hi}     & 55.89 & \textbf{11.70} & 270.3 \\
\bottomrule
\end{tabular}
\end{center}

\noindent
\textbf{Rank by feature risk:} refined, good, karras, forced, cluster. \\
\textbf{Rank by FID:}\hphantom{xxxxxx} good, refined, karras, forced, cluster.

\noindent The feature risk reproduces the FID order in $3$ of $5$
positions exactly, with only the top two (\texttt{good} vs
\texttt{refined}, both low-FID) swapped --- within the noise of a
single-batch, leading-order, $28$-node estimate of $g$. Decisively, it
fixes the two failures of the pixel certificate:
\begin{itemize}
\item \texttt{good} vs \texttt{forced\_smin}: pixel certificate
$26.0\approx26.4$ (indistinguishable); feature risk $1.46$ vs $4.00$
(correctly separated, matching FID $13$ vs $88$).
\item \texttt{cluster\_hi}: pixel certificate middling ($55.9$); feature
risk correctly worst ($11.70$, matching FID $270$).
\end{itemize}

\section{Consequences}
\label{sec:consequences}

\begin{itemize}
\item \textbf{The KL/$L^2$ certificate (PDF~1) is correct for what it
bounds} --- NLL / pathwise-$L^2$ risk --- and is not an FID certificate.
This is consistent with the broader literature: KL/ELBO-optimal schedules
are known to differ from FID-optimal ones, which is why schedule-design
work (EDM, AYS) tunes to FID rather than minimizing a likelihood bound.
\item \textbf{The locked empirical wins are unaffected.} They are measured
Clean-FID results; their validity does not depend on certificate
faithfulness. The certificate supplies the \emph{form}
$m^\star\propto d_s^{1/(p+1)}$ and the $k{=}2$ weight tunes it toward FID
empirically.
\item \textbf{The FID-faithful certificate is the feature-space
Wasserstein risk} $Q^{\varphi}_s$ with $d_s^{\varphi}=d_s\,g$, validated
above. It replaces the tuned scalar $k$ with the measured profile
$g(\sigma)$, making the FID claim certified rather than tuned.
\end{itemize}

\section{Caveats and next steps}
\label{sec:caveats}

\begin{itemize}
\item The validation is leading-order (covariance term of
$\mathrm{FID}_{\mathrm{disc}}$; the mean-shift/bias term is dropped) and
uses a single calibration batch with a $28$-node $g$ profile and a
differentiable-bicubic resize proxy for clean-FID's exact resizer. The
rank agreement is therefore rank-level, not absolute. A $5$-seed,
finer-$g$, exact-resize study is the confirmation run.
\item The bias term $\lVert\mathbb E[\Delta\varphi]\rVert^2$ should be
added for grids with systematic (non-zero-mean) error; the present
validation captures the covariance/variance term.
\item The end-to-end test is to \emph{optimize}
$m^\star_{\mathrm{FID}}\propto(d_s g)^{1/(p+1)}$ and measure whether it
beats the $k{=}2$-tuned pointwise grid on locked Clean-FID. This converts
the FID claim from empirically-tuned to certified.
\end{itemize}

\end{document}
